\section{Why the Law Values Celibacy without Making It the Essence of Orders}

\subsection{Christological configuration}

The priest acts sacramentally in the person of Christ the Head and Shepherd. Christ's own virginal, wholly filial life makes celibacy a fitting personal sign of the minister's total availability to the Father and the people entrusted to him. The sign is powerful because it is embodied, costly, and free. It would lose its theological character if imposed without discernment, treated as a guarantee of maturity, or reduced to labor availability.

The argument is one of fittingness, not sacramental necessity. A married presbyter is no less configured to Christ by ordination. Vatican II says this directly in \emph{Presbyterorum ordinis} 16, pointing to the primitive Church and the continuing Eastern discipline. Catholic theology therefore holds together a high account of celibacy and full recognition of married priesthood.

\subsection{Ecclesial and nuptial meaning}

Celibacy can express pastoral charity: the cleric gives himself to a local Church and anticipates the complete communion of the kingdom. John Paul II's \emph{Pastores dabo vobis} 29 describes the commitment as a gift of self rather than a merely negative restriction. Paul VI's \emph{Sacerdotalis caelibatus} likewise sets out Christological, ecclesial, pastoral, and eschatological dimensions.

Nuptial language must be disciplined. The Church is the Bride of Christ, and an ordained minister can represent Christ's pastoral self-gift. He does not thereby acquire a sacramental marriage to a parish that cancels the real marriage of an Eastern priest or permanent deacon. Nor does ministry give a cleric ownership of the community or exemption from ordinary boundaries.

\subsection{Eschatological witness}

Celibacy witnesses to the resurrection, where human communion is not structured by contracting marriage (Matt 22:30), and to the kingdom already breaking into the present age. Marriage also witnesses to the covenant of Christ and the Church (Eph 5:25--32). Catholic tradition needs both signs. They illuminate rather than invalidate one another.

\subsection{Freedom, charism, and ecclesiastical selection}

No man has a right to ordination. The Church may select candidates according to a discipline that requires the demonstrated reception of a celibate charism. The candidate's freedom is real but situated: formation must permit him to decline without shame, investigate affective maturity, and prevent a desire for ministry from disguising an inability to embrace the obligation. Ecclesial authority, for its part, cannot manufacture the charism merely by imposing a rule.

The Latin requirement is ecclesiastical law and can be modified by the competent authority. Calling it mutable does not mean it is dispensable by a candidate, a bishop acting beyond his competence, or private agreement. Conversely, treating its theological fittingness as important does not convert it into a divine-law condition of valid ordination.

\subsection{Chastity is not a safeguarding system}

Celibacy neither causes abuse by definition nor prevents it by itself. Sexual coercion and abuse involve violations of persons, power, office, and sometimes criminal law as well as chastity. Safe ministry requires screening, formation, supervision, reporting systems, civil-law compliance, canonical investigation, victim-centered care, and accountability. A spiritual vocabulary that calls abuse only a ``failure of celibacy'' obscures the harmed person and the abuse of authority. A policy vocabulary that treats celibacy as the single causal variable is equally inadequate.

\begin{fourstable}{Claim}{True element}{Required limit}{Canonical significance}
Celibacy is a gift & It can embody undivided pastoral charity and the kingdom. & The gift requires discernment and does not prove personal holiness. & The Church may require it for admission without declaring marriage inferior.\\
Celibacy is discipline & Its Latin juridic requirement is ecclesiastical and has admitted exceptions. & Its theological meanings and the moral obligation of chastity are not arbitrary policy. & Change belongs to competent authority, not private nonobservance.\\
Married clergy are traditional & Married ministry is apostolic-era and continuously legitimate in the East. & Eastern law is not simply the Latin rule plus a dispensation. & Each Church's common and particular law must be applied on its own terms.\\
Continence has ancient roots & Western texts clearly legislate it by the fourth century and claim apostolic warrant. & Surviving evidence does not establish uniform first-century enforcement beyond dispute. & Present law does not depend on resolving the historical controversy.\\
\end{fourstable}
